% ================================================================
\section{Convergence Analysis}\label{sec:convergence}

We now formalize the two key convergence properties that
differential flatness enables: (i)~a quadratic gain
approximation error and (ii)~exponential time-contraction
of the closed-loop system.

\begin{proposition}[Quadratic gain approximation]\label{prop:gain_approx}
Let $K(\balpha)$ denote the DARE gain~\eqref{eq:kinf} at
flat parameter $\balpha$, and let
$\tilde{K}(\balpha) = K_0 + \sum_{i=1}^{4}
(\alpha_i - \alpha_{0,i}) \, \partial_i K$
be the first-order approximation~\eqref{eq:K_interp}.
Suppose $K(\cdot)$ is twice continuously differentiable
in a neighborhood of $\balpha_0$
(guaranteed when $(A_d(\balpha), B_d(\balpha))$ is
stabilizable and $(Q^{1/2}, A_d(\balpha))$ is
detectable~\cite{konstantinov1993}).
Then
\begin{equation}\label{eq:gain_quad}
  \|K(\balpha) - \tilde{K}(\balpha)\|_F
  \;\leq\; \tfrac{1}{2} \, L_K \, \|\balpha - \balpha_0\|^2,
\end{equation}
where $L_K = \max_i \sup_{\balpha} \|\partial^2 K / \partial \alpha_i \partial \alpha_j\|_F$
is the Lipschitz constant of the gain Jacobian.
\end{proposition}

\begin{proof}
This is the standard second-order Taylor remainder applied
to the matrix-valued map $\balpha \mapsto K(\balpha) \in \R^{m \times n}$.
The twice-differentiability of $K$ with respect to
the system parameters $(A, B)$ follows from the implicit
function theorem applied to the DARE~\cite{hewer1971,konstantinov1993},
composed with the smooth (trigonometric) dependence of
$(A_d(\balpha), B_d(\balpha))$ on $\balpha$
from~\eqref{eq:Ac}--\eqref{eq:disc}.
\end{proof}

\textit{Why flatness matters.}
Without differential flatness, the scheduling variable
would be the full $n$-dimensional state, and the gain
Hessian $\partial^2 K / \partial x_i \partial x_j$ would
be an $O(n^2)$-parameter tensor.
Flatness compresses the scheduling space from $\R^n$ to
$\R^4$, so the Hessian has only $\binom{4+1}{2} = 10$
independent components---making~\eqref{eq:gain_quad}
a tight, low-dimensional bound.

\begin{proposition}[Exponential time-contraction]\label{prop:contraction}
Let $A_{\mathrm{cl}}(\balpha) = A_d(\balpha) - B_d(\balpha) K(\balpha)$
be the closed-loop matrix at flat parameter $\balpha$.
The DARE solution guarantees
$\bar\rho(A_{\mathrm{cl}}) < 1$
(all eigenvalues strictly inside the unit circle),
so the error state contracts exponentially:
\begin{equation}\label{eq:contraction}
  \|\delta x_k\| \;\leq\; c \; \bar\rho(A_{\mathrm{cl}})^k \; \|\delta x_0\|,
  \qquad c \geq 1,
\end{equation}
where $\delta x_k = x_k - x^*(\balpha)$,
$\bar\rho(A_{\mathrm{cl}}) < 1$ is the spectral radius,
and $c$ depends on the condition number of the eigenvector
matrix of $A_{\mathrm{cl}}$.
\end{proposition}

\textit{Contraction rate vs.\ constraint conservatism.}
For the barrier-augmented variant (Section~\ref{sec:barrier}),
$R_\mu = R + H_\mu(\balpha)$ increases the effective input
cost, producing a less aggressive gain $K_\mu$ and a larger
spectral radius $\bar\rho(A_{\mathrm{cl}}^\mu) > \bar\rho(A_{\mathrm{cl}})$.
Thus the barrier parameter $\mu$ explicitly trades contraction
speed for constraint margin: larger $\mu$ yields slower
convergence but greater safety distance from actuator limits.

\begin{corollary}[Tracking error bound]\label{cor:tracking}
When the first-order gain approximation $\tilde{K}(\balpha)$
is used instead of the exact gain $K(\balpha)$,
the closed-loop matrix becomes
$\tilde{A}_{\mathrm{cl}} = A_d - B_d \tilde{K}$.
By~\eqref{eq:gain_quad}, the perturbation
$\Delta K = K - \tilde{K}$ satisfies
$\|\Delta K\| \leq \frac{1}{2} L_K \|\delta\balpha\|^2$.
Standard perturbation theory for discrete-time
systems~\cite{konstantinov1993} then gives
\begin{equation}\label{eq:tracking_bound}
  \bar\rho(\tilde{A}_{\mathrm{cl}})
  \;\leq\; \bar\rho(A_{\mathrm{cl}})
  + \|B_d\| \cdot \tfrac{1}{2} L_K \|\delta\balpha\|^2
  + O(\|\delta\balpha\|^4).
\end{equation}
The approximate closed loop remains stable whenever
$\|\delta\balpha\|$ is small enough that the right-hand
side is less than~$1$.
For the quadrotor with $\|\delta\balpha\| \leq 3$\,m/s$^2$,
our numerical tests (Section~\ref{sec:experiments})
confirm mean gain error of $4\%$ and stable tracking
across 100 random $\balpha$ samples.
\end{corollary}
